\documentclass[10pt]{article}

\usepackage[utf8]{inputenc}

\usepackage[T1]{fontenc}

\usepackage{amsmath}

\usepackage{amsfonts}

\usepackage{amssymb}

\usepackage[version=4]{mhchem}

\usepackage{stmaryrd}

\usepackage{bbold}

\usepackage{graphicx}

\usepackage[export]{adjustbox}

\graphicspath{ {./images/} }

\usepackage{mathrsfs}



\begin{document}

гич. произведения $E_{1} \times E_{2}$ в $E_{1} \bigotimes E_{2}$. П р о е к т и вн о й (соответственно, и в д у к т и в н о й) топологией в $E_{1} \otimes E_{2}$ наз. самая сильная среди всех локально выпуклых топологий в $E_{1} \otimes E_{2}$, для к-рых отображение $b$ непрерывно (соответственно, раздельно непрерывно). Хотя эта терминология не вполне последовательна, она является общепринятой. Л.ю. п., получающееся в результате наделения векторного пространства $E_{1} \bigotimes E_{2}$ проективной (соответственно, индуктивной) топологией, обозначается символом $E_{1} \bigotimes_{\pi} E_{2}\left(E_{1} \otimes_{i} E_{2}\right)$, а его пополнение - символом $E_{1} \widehat{\bar{\otimes}}_{\pi} E_{2}\left(E_{1} \widehat{\widehat{\otimes}}_{i} E_{2}\right)$. Пространства $E_{1} \otimes_{\pi} E_{2}$ и $E_{1} \otimes_{i} E_{2}$ наз. ло к а ль н о в ы п у клыми тензорными произведения и иоответствующих л. в. п., а их пополнения - полными локально выпуклыми тензорными произведениями. Существуют - помимо введенных - и другие локально выпуклые тензорные произведения; они получаются путем наделения алгебраического тензорного произведения топологиями, отличными от описанных. Многие свойства тензорных произведений упрощаются, когда одно из пространств-сомножителей - ядерное пространство.



П р и м е р ы. Л. в. п. $S(\mathbb{R} n) \widehat{\widehat{\otimes}} \pi S(\mathbb{R} k), S\left(\mathbb{R}^{n}\right) \hat{\bigotimes}_{i} S\left(\mathbb{R}^{k}\right)$ и $S(\mathbb{R} n+k)$ канонически изоморфны (изоморфизм первых двух л. в. П. является следствием того, что всякое раздельно непрерывное билинейное отображение произведения пространств Фреше в произвольное л. в. п. непрерывно). Л. в. П. $D(\mathbb{R} n) \widehat{\otimes}{ }_{i} D(\mathbb{R} k)$ и $D(\mathbb{R} n+k)$ также канонически изоморфны; векторные пространства $D\left(\mathbb{R}^{n}\right) \hat{\bigotimes}_{\pi} D(\mathbb{R} k)$ и $D\left(\mathbb{R}^{n+k}\right)$ канонически изоморфны, хотя топологии в них не совпадают [8], [9].



Двойственность. Важную роль при исследовании л. в. П. играет использование двойственности между л. в. п. и его сопряженным. В частности, нек-рые свойства л. в. п. зависят лишь от объема сопряженного пространства. Так, если $E$ - л. в. п., и $E^{\prime}$ - его сопряженное, то для всех локально выпуклых топологий в $E$, coгласующихся с двойственностью между $E$ и $E^{\prime}$, ограниченные множества - одни и те же и замкнутые выпуклые множества также одни и те же.



Теория двойственности оказывается полезной при исследовании полных пространств. Так, л. в. п. (соответственно, метризуемое л. в. п.) $E$ полно в том и только в том случае, если всякое гиперподпространство (соответственно, выпуклое подмножество) его сопряженного $E^{\prime}$, обладающее замкнутыми в топологии $\sigma\left(E^{\prime}, E\right)$ пересечениями с полярами всех окрестностей нуля пространства $E$, само замкнуто в этой топологии (т е о р м Б К р е й на - ШІІ у лья н а).



В связи с этим вводятся следующие определения. Л. в. Ш. наз. $B_{r^{-}}$П о л н ы м (соответственно, с о в е рші енно полны с т р а нс т в м всякое всюду плотное линейное подпространство (соответственно, линейное подпространство, абсолютно выпуклое подмножество, выпуклое подмножество) пространства $\left(E^{\prime}, \sigma\left(E^{\prime}, E\right)\right)$, обладающее замкнутыми пересечениями с полярами всех окрестностей нуля из $E$, само является замкнутым. Эти классы пространств играют важную роль в обобщениях теорем Банаха о замкнутом графике и открытом отображении (см. ниже). Полные метризуемые л. в. п. и сильные сопряженные к рефлексивным (см. ниже) метризуемым л. в. п. входят в каждый из этих классов; в то же время пространства $D$ и $D^{\prime}$ не входят ни в один из них. Классы гиперполных пространств и пространств Крейна - Шіульяна не совпадают, однако до сих пор (1984) не известно, совпадают или нет классы $B_{r}$-полных и гишерполных пространств. C помоцью теории двойственности доказываются также следующие предложения о компактных подмножествах л. в. П. 1) Пусть $E$ - л. в. п. и $H$ - подмножество $E$, замкнутая выпуклая оболочка которого полна в топологии Макки. Если всякая последовательность әлементов из $H$ обладает в $E$ предельной точкой, то $H$ относительно компактно (т е о р е м а Әб е р л е й н а). 2) Пусть $E$ - метризуемое л. в. п., $\left\{x_{n}\right\}$ - последовательность в $E$, всякая подпоследовательность к-рой обладает в $\left(E, \sigma\left(E, E^{\prime}\right)\right)$ предельной точкой. Тогда из последовательности $\left\{x_{n}\right\}$ можно извлечь сходящуюся подпоследовательность (т е о р е м а III м у л ь я н а). 3) Пусть $B$ - компактное подмножество в отделимом л. в. П. $E$ и $C$ - замкнутая выпуклая оболочка $B$. Тогда $C$ компактно в том и только в том случае, когда оно полно в топологии Макки (т е о р е м а $К$ р е й н а).



Л. в. п. $E$ наз. п ол у р е ф ле к с и вны м (соответственно, р е ф л е к с и в н ы м), если канонич. вложение $x \mapsto[g \mapsto g(x)], E \rightarrow\left(E_{\beta}^{\prime}\right)_{\beta}^{\prime}$ представляет собой изоморфизм векторных пространств (соответственно, изоморфизм Т. в. П.). Для того чтобы л. в. п. $E$ было полурефлексивно, необходимо и достаточно, чтобы всякое его ограниченное подмножество было относительно компактно в тонологии $\sigma\left(E, E^{\prime}\right)$; для того чтобы оно было рефлексивным, необходимо и достаточно, чтобы оно было полурефлексивным и бочечным пространством.



Отображения Т. в. п. 1. Т е о р е м ы о з а м к н утом графике и о ткрытом отображен и и. Линейное отображение Т. в. п. $E_{1}$ в Т. в. п. $E_{2}$ наз. то по ло г и ч е ским гомом о ро и зм о м, если оно переводит всякое открытое в $E_{1}$ множество в множество, открытое в $f\left(E_{1}\right)$ (в топологии, индуцированной топологией пространства $\left.E_{2}\right) . \Gamma \mathrm{p}$ аф и к о отображения $f: E_{1} \rightarrow E_{2}$ наз. множество $\left\{\left(x, f(x): x \in E_{1}\right\} \subset E_{1} \times E_{2}\right.$.



\begin{center}

\includegraphics[max width=\textwidth]{2023_07_09_d167bf39114f6ebd8cb6g-1(2)}

\end{center}



\includegraphics[max width=\textwidth, center]{2023_07_09_d167bf39114f6ebd8cb6g-1}

графике (соответственно, о гомоморфизме), если для всех $E_{1} \in \mathscr{O}_{1}, E_{2} \in \mathscr{O}_{2}$ всякое линейное отображение $f: E_{1} \rightarrow E_{2}$, обладающее замкнутым в Т. в. П. $E_{1} \times E_{2}$ графиком, непрерывно (соответственно, если всякос непрерывное линейное отображение $E_{2}$ на $E_{1}$ являетсл топологич. гомоморфизмом). Если $\mathscr{O}-$ класс всех полных метризуемых Т. в. п., то для пары (£, справедливы как теорема о замкнутом графике, так и теорема о гомоморфизме (т е о р е м а Б а н а х а). Этот результат усилен: пусть $\mathscr{O}_{1}-$ класс всех отделимых л. в. п., являющихся индуктивными пределами семейств банаховых пространств, а $\mathscr{O}_{2}-$ наименьший класс л. в. п., содержащий все полные метризуемые л. в. П. и замкнутый относительно образования проективных и индуктивных пределов счетных семейств входящих в него пространств. Тогда для пары $\left(\mathscr{O}_{1}, \mathscr{O}_{2}\right)$ справедливы теоремы о замкнутом графике и гомоморфизме (т е о р е м а $\mathrm{P}$ а й к о в а) [все встреqающиеся в функциональном анализе сепарабельные л. в. п. в их неослабленной топологии входят сразу в оба эти класca]. Фактически сформулированное утверждение доказано для нек-рого, более широкого, чем $\mathscr{O}_{2}$, класса $\mathscr{O}_{2}^{\prime}$ Т. в. п. и для многозначных линейных отображений; описан [7] еще один класс $\mathfrak{O}_{2}^{\prime \prime} \supset \mathcal{O}_{2}^{\prime}$ Т. в. п., к-рый может играть в сформулированном утверждении роль класса $\mathscr{O}_{2}$ - т. н. П р о с т р а н с т в а с с е т ь ю.



Пусть $\mathscr{O}, \mathscr{O}_{\mathscr{P}}, \mathscr{O}_{\mathscr{P}}^{r}$, соответственно,- классы всех бочечных, всех совершенно полных и всех $B_{r}$-полных



\includegraphics[max width=\textwidth, center]{2023_07_09_d167bf39114f6ebd8cb6g-1(1)}

рема о замкнутом графике, а для пары $\left(\mathscr{O}_{\mathscr{O}}^{\mathscr{O}} \mathscr{\mathcal { P }}\right)-$ теорема о гомоморфизме.





\end{document}